% ════════════════════════════════════════════
% 情景设定（5.4.2）
% ════════════════════════════════════════════
\subsubsection{情景设定}

本节不重新估计模型，而是读取回归阶段保存的最终模型与 500 条清洗数据，构造反事实情景。为使不同处理组之间只相差一个变量，设定代表性画像为 Photo、主题 1、7 月、周四、上午、粉丝规模取样本中位数 129,600 人。该组合的各个水平都在样本中有充足观测，避免外推到稀疏单元。付费情景只把付费标记改为 1，自然情景改为 0，其余分量保持不变。需要说明的是，粉丝规模与月份在样本中近乎一一对应，月份对粉丝规模对数的解释度达 0.9917，M2 以月份固定效应吸收规模趋势；因此画像中的粉丝规模取值用于标明情景所处的发展阶段，在样本的规模—月份对应关系下，固定月份、内容类型与时段后，处理组与对照组之间即满足“粉丝量相同”的口径。

除代表性画像外，另做全样本反事实。对第 $i$ 条帖子，保留其全部协变量，分别把付费标记替换为 1 与 0，得到两个潜在结果；前者表示该帖子若被付费推广时的曝光，后者表示保持自然状态时的曝光。内容类型与内容主题的反事实同理，只替换对应的哑变量。这样构造的样本平均处理效应既利用了全样本的协变量分布，又不受各处理组样本量不平衡的影响。

在得到反事实预测后，需要选择把对数尺度结果还原到曝光次数并量化不确定性的方法。针对这一任务，需在系数直接指数化、Duan 烟雾还原、参数模拟与非参数 Bootstrap 之间择优，从偏差修正、不确定性刻画与计算成本三个维度对比其适用性，如表\ref{tab:algo4}所示。

\begin{longtable}{>{\centering\arraybackslash}p{0.18\textwidth}>{\centering\arraybackslash}p{0.20\textwidth}>{\centering\arraybackslash}p{0.28\textwidth}>{\centering\arraybackslash}p{0.22\textwidth}}
  \caption{效应测算方法对比}
  \label{tab:algo4} \\
  \toprule
  方法 & 偏差修正 & 不确定性刻画 & 计算成本 \\
  \midrule
  \endfirsthead
  \caption{效应测算方法对比（续）} \\
  \toprule
  方法 & 偏差修正 & 不确定性刻画 & 计算成本 \\
  \midrule
  \endhead
  \bottomrule
  \endfoot
  系数直接指数化 & 只给出几何均值倍数 & 用系数区间近似 & 最低 \\
  Duan 烟雾还原 & 非参数修正均值 & 需与模拟或重抽样结合 & 低 \\
  参数模拟 & 沿用烟雾还原 & 传播 HC3 协方差 & 中等 \\
  非参数 Bootstrap & 沿用烟雾还原 & 不依赖渐近正态 & 较高 \\
\end{longtable}

如表\ref{tab:algo4}所示，直接对系数取指数只能得到几何均值意义上的倍数，会低估曝光次数的算术均值；烟雾还原用残差的经验分布修正这一偏差，且不额外假定对数正态\cite{duan1983}。不确定性方面，参数模拟从 HC3 协方差中抽取系数向量，能同时传播点估计与协方差信息，计算量适中；非参数 Bootstrap 虽更稳健，但需要重复重估模型，在当前样本量下收益有限。综合以上分析，本部分选择烟雾还原结合参数模拟的方案，Bootstrap 仅作为方法学上的备选。
